\documentclass[11pt]{article}

\usepackage[letterpaper,margin=0.82in,headheight=14pt]{geometry}
\usepackage[T1]{fontenc}
\usepackage[utf8]{inputenc}
\usepackage{lmodern}
\usepackage{microtype}
\usepackage{amsmath,amssymb,mathtools,amsthm}
\usepackage{booktabs,longtable,tabularx,array}
\usepackage{enumitem}
\usepackage{xcolor}
\usepackage{xurl}
\usepackage{hyperref}
\usepackage{fancyhdr}
\usepackage{titlesec}
\usepackage[most]{tcolorbox}
\usepackage{listings}
\usepackage{lastpage}

\definecolor{navy}{HTML}{17324D}
\definecolor{teal}{HTML}{087E8B}
\definecolor{crimson}{HTML}{A3323D}
\definecolor{amber}{HTML}{B06A00}
\definecolor{ink}{HTML}{202A33}
\definecolor{slate}{HTML}{5D6B78}
\definecolor{mist}{HTML}{F1F5F7}
\definecolor{palegreen}{HTML}{EAF5F1}
\definecolor{palered}{HTML}{FAECEE}
\definecolor{paleyellow}{HTML}{FFF6E5}

\hypersetup{
  colorlinks=true,
  linkcolor=navy,
  citecolor=teal,
  urlcolor=teal,
  pdftitle={Cold Review of the Current ALCI-RCC5 Development},
  pdfauthor={Independent technical review},
  pdfsubject={Overview-paper audit, Lean artifact audit, and full-logic decidability analysis},
  pdfkeywords={ALCI, RCC5, decidability, Lean, description logic, cold review}
}

\titleformat{\section}{\Large\bfseries\color{navy}}{\thesection}{0.65em}{}
\titleformat{\subsection}{\large\bfseries\color{navy}}{\thesubsection}{0.55em}{}
\titleformat{\subsubsection}{\normalsize\bfseries\color{ink}}{\thesubsubsection}{0.5em}{}
\titlespacing*{\section}{0pt}{2.0ex plus .4ex minus .2ex}{0.8ex}
\titlespacing*{\subsection}{0pt}{1.5ex plus .3ex minus .2ex}{0.5ex}

\pagestyle{fancy}
\fancyhf{}
\lhead{Cold review: ALCI\(_{\mathrm{RCC5}}\)}
\rhead{9 September 2026}
\cfoot{\thepage\ / \pageref*{LastPage}}
\renewcommand{\headrulewidth}{0.35pt}

\setlist[itemize]{leftmargin=1.45em,itemsep=2.3pt,topsep=3pt,parsep=1pt}
\setlist[enumerate]{leftmargin=1.65em,itemsep=2.3pt,topsep=3pt,parsep=1pt}
\renewcommand{\arraystretch}{1.16}
\emergencystretch=2em
\newcolumntype{L}[1]{>{\raggedright\arraybackslash}p{#1}}
\newcolumntype{Y}{>{\raggedright\arraybackslash}X}

\newtcolorbox{verdictbox}[1][]{
  enhanced,breakable,colback=palered,colframe=crimson,
  boxrule=0.8pt,arc=2pt,left=8pt,right=8pt,top=7pt,bottom=7pt,
  title={#1},fonttitle=\bfseries,before upper=\raggedright
}
\newtcolorbox{resultbox}[1][]{
  enhanced,breakable,colback=palegreen,colframe=teal,
  boxrule=0.7pt,arc=2pt,left=8pt,right=8pt,top=7pt,bottom=7pt,
  title={#1},fonttitle=\bfseries,before upper=\raggedright
}
\newtcolorbox{cautionbox}[1][]{
  enhanced,breakable,colback=paleyellow,colframe=amber,
  boxrule=0.7pt,arc=2pt,left=8pt,right=8pt,top=7pt,bottom=7pt,
  title={#1},fonttitle=\bfseries,before upper=\raggedright
}

\lstdefinelanguage{Lean}{
  morekeywords={theorem,def,structure,where,Prop,Type,forall,exists,by,if,then,else,match,with},
  morecomment=[l]{--},sensitive=true
}
\lstset{
  language=Lean,basicstyle=\ttfamily\footnotesize,breaklines=true,
  columns=fullflexible,keepspaces=true,frame=single,framerule=0.3pt,
  rulecolor=\color{slate!50},backgroundcolor=\color{mist},
  xleftmargin=0.4em,xrightmargin=0.4em
}

\newcommand{\Logic}{\ensuremath{\mathcal{ALCI}_{\mathrm{RCC5}}}}
\newcommand{\PO}{\ensuremath{\mathsf{PO}}}
\newcommand{\PP}{\ensuremath{\mathsf{PP}}}
\newcommand{\PPI}{\ensuremath{\mathsf{PPI}}}
\newcommand{\DR}{\ensuremath{\mathsf{DR}}}
\newcommand{\EQ}{\ensuremath{\mathsf{EQ}}}
\newcommand{\sat}{\mathsf{Sat}}
\newcommand{\unsat}{\mathsf{Unsat}}
\newcommand{\code}[1]{\texttt{#1}}
\newcommand{\sevmajor}{\textcolor{crimson}{\textbf{Major}}}
\newcommand{\sevmod}{\textcolor{amber}{\textbf{Moderate}}}
\newcommand{\sevminor}{\textcolor{slate}{\textbf{Minor}}}

\begin{document}
\pagenumbering{roman}

\begin{titlepage}
  \color{ink}
  \vspace*{0.5in}
  \noindent{\color{teal}\rule{\linewidth}{2.4pt}}\\[0.35in]
  {\Huge\bfseries\color{navy} Cold Review of the Current\\[0.12in]
  \Logic\ Development}\par
  \vspace{0.22in}
  {\Large Overview-paper audit, Lean artifact review,\\
  and a renewed attack on full decidability}\par
  \vspace{0.34in}
  {\large Review date: 9 September 2026}\par
  \vspace{0.5in}

  \begin{tcolorbox}[colback=mist,colframe=navy,boxrule=0.9pt,arc=2pt,
    before upper=\raggedright]
  \textbf{Submission reviewed.} Attached archive \code{alcircc5-master.zip},
  archive comment / commit identifier
  \code{795b2707b72e927c56b465ef0413b77aeb9a66a6}; overview paper
  \code{papers/overview\_arxiv.tex} and the five canonical files under
  \code{formal/}.
  \end{tcolorbox}

  \vfill
  {\large Independent adversarial adequacy review}\par
  \vspace{0.08in}
  {\small Claims are graded by what the archived sources actually establish,
  not by filenames, comments, or intended future bridges.}\par
  \vspace{0.25in}
  \noindent{\color{teal}\rule{\linewidth}{2.4pt}}
\end{titlepage}
\clearpage
\pagenumbering{arabic}

\begin{abstract}
This report begins with the 52-page overview paper and then audits the current
Lean sources, executable probes, and the proposed routes to full
\Logic{} decidability.  The central verdict is deliberately two-part.  First,
the archive gives strong source-level evidence for a genuine, executable,
reportedly kernel-checked decision procedure for \emph{bare concept
satisfiability} in the
polarity-correct \(\forall\PO\)-free fragment, including the bridge to arbitrary
nonempty-set semantics.  Second, it does \emph{not} prove decidability of the
full logic.
The current cone construction is complete in the no-false-negative direction
for the full language, but its model-extraction/soundness argument uses the
fragment restriction exactly at positive \(\forall\PO\).

The renewed attack did not resolve the problem left open by the primary 2006
comparator and unresolved in this snapshot.  It did produce a sharper
diagnosis: explicit unsatisfiable formulas survive the current unary-signature
architecture because non-singleton RCC5 compositions impose obligations
between independently generated witnesses.  A converse satisfiable example
shows that declaring all such cross-pairs to be \PO{} is not a repair.  The
report also gives (i) a witness-generated countable-model lemma, and (ii) a
two-layer encoding of ordered-disjoint frames into one strict partial order.
The latter absorbs RCC5 heredity into transitivity, but introduces a synchronized
fixed-point-free involution not covered by the standard two-variable/one-
transitive-relation theorem.  These are useful reductions, not a decision proof.

For publication, the recommendation is major rewrite: extract a focused paper
on the reported fragment theorem and normal form, and recast the full-logic discussion
as a research program with every missing implication stated explicitly.
\end{abstract}

\begingroup
\setcounter{tocdepth}{1}
\small
\tableofcontents
\endgroup
\newpage

\section{Scope, materials, and review standard}

\subsection{Snapshot and objects reviewed}

The attached ZIP, rather than an assumed later repository state, is the review
object.  Its SHA-256 digest is
\begin{center}
\small\code{4a0415b3e3202595d1b1290e6cc741068d0423c0c2877e815b9b5fcbe96bd2a8}.
\end{center}
The ZIP comment identifies commit
\code{795b2707b72e927c56b465ef0413b77aeb9a66a6}.  The overview source digest is
\code{be28d3da\ldots429e}; its included PDF has 52 pages and digest
\code{e684c506\ldots322a}.  Full hashes of the formal sources appear in
Appendix~\ref{app:repro}.

I attempted to compare the public GitHub head and to obtain a fresh Lean
toolchain.  Network access to those endpoints was denied in this execution
environment, and no \code{lean}, \code{elan}, or \code{lake} executable was
preinstalled.  Accordingly, this is a cold review of the identified archive,
not a statement about commits after it, and the Lean judgment below separates
source audit from independent compilation.

\subsection{Four evidential levels}

I use four labels throughout:
\begin{enumerate}
  \item \textbf{Source-audited Lean declaration}: a declaration present in a
        canonical Lean source, with no project-specific axiom or proof escape
        visible in the source.  Without a local build this is not itself a
        kernel-verification claim.
  \item \textbf{Archived build claim}: a build/axiom result recorded in the
        submission's transcripts, but not reproduced in this environment.
  \item \textbf{Executable evidence}: a Python checker or finite exhaustive
        test reproduced during this review.  This is regression evidence, not
        a proof of an unbounded theorem.
  \item \textbf{Mathematical/prose claim}: a paper argument, conjecture, or
        proposed bridge that is not an instantiated Lean theorem.
\end{enumerate}
This discipline matters here: the repository contains excellent source-audited,
reportedly kernel-checked combinators whose decisive completeness hypotheses
are still structure fields or theorem parameters.

\section{Executive verdict}

\begin{verdictbox}[Overall verdict]
\textbf{Full \Logic{} decidability is not proved and was not cracked in this
review.}  The problem is unresolved in the submitted snapshot and inspected
corpus.  The
\(\forall\PO\)-free concept-satisfiability theorem is the strongest completed
result and appears technically real.  The overview paper, however, currently
promotes an architecture-specific conjectural route (F6 plus a uniformization
bridge) into a ``one keystone'' reduction.  As a formal-results paper it
requires a major rewrite; as two papers---a focused fragment/formalization paper
and a candid research-program/history paper---the material is much stronger.
\end{verdictbox}

\begin{table}[ht]
\centering
\small
\begin{tabularx}{\textwidth}{L{0.27\textwidth}L{0.18\textwidth}Y}
\toprule
Claim & Verdict & Basis / qualification \\
\midrule
Full concept satisfiability is decidable & \textcolor{crimson}{\textbf{Unresolved here}} &
No complete finite-presentation theorem; current cone extraction fails with
positive \(\forall\PO\). \\
Full unsatisfiability has a sound cone test & \textcolor{teal}{\textbf{Supported}} &
\code{coneScheme\_complete} gives full-language no-false-negatives; the
contrapositive theorem is explicit. \\
\(\forall\PO\)-free concept satisfiability is decidable &
\textcolor{teal}{\textbf{Strongly supported}} & End-to-end source chain from
raw polarity test through arbitrary nonempty-set semantics; clean external rebuild still
needed. \\
Normal form / set representation & \textcolor{teal}{\textbf{Supported}} &
Substantive Lean development, but the overview cites the wrong file for the
converse assembly. \\
F6 is the sole remaining theorem & \textcolor{crimson}{\textbf{Not established}} &
F6 is informal and the F6\(\Rightarrow\)uniformization\(\Rightarrow\)finite-
checker-completeness bridge is open. \\
Overview ready as a formal-results paper & \textcolor{crimson}{\textbf{No}} &
Major claim calibration, status consolidation, scope, citation, and
reproducibility changes are required. \\
\bottomrule
\end{tabularx}
\caption{Current-state verdict at archive commit \code{795b2707\ldots66a6}.}
\end{table}

\paragraph{Confidence calibration.}
I found no evidence that Lean accepted a false theorem.  The problem is almost
the opposite: several prose claims outrun the hypotheses of sound, honest Lean
theorems.  The fragment result deserves a high-confidence label subject to a
fresh pinned build.  Any headline resolution of the full logic deserves a
low-confidence label until a model-existence/completeness theorem is supplied.

\section{Cold review of the overview paper}

\subsection{What is strong}

The paper has real strengths worth preserving:
\begin{itemize}
  \item It records the long sequence of failed proof routes more candidly than
        most research narratives and often distinguishes abstract relation-
        algebra semantics from topological semantics correctly.
  \item The ordered-disjoint normal form is a useful conceptual compression:
        \PP{} becomes strict order, \DR{} becomes hereditary conflict, and
        \PO{} is the residual incomparable non-conflict relation.
  \item The final cone-scheme section points to a substantial Lean theorem
        reported to build
        for the whole polarity-correct \(\forall\PO\)-free fragment, not merely
        the older three ``pure quadrants.''
  \item The paper exposes counterexamples and retractions instead of deleting
        them from the historical record.
  \item The abstract-to-set satisfiability bridge and the explicit theorem-name
        map are exactly the right reproducibility instincts, even though the
        map needs correction.
\end{itemize}

\subsection{Ranked findings}

\small
\begin{longtable}{L{0.055\textwidth}L{0.115\textwidth}L{0.37\textwidth}L{0.36\textwidth}}
\toprule
ID & Severity & Finding & Required correction \\
\midrule
\endfirsthead
\toprule
ID & Severity & Finding & Required correction \\
\midrule
\endhead
O1 & \sevmajor & ``Reducing it to one keystone'' is not a proved reduction.
F6 is informal, and the finite-code completeness theorem is not derived from it.
& Retitle and display every missing arrow; call F6 a candidate sufficient route
within one architecture. \\
O2 & \sevmajor & The identity-selector obstruction is described as an iff /
theorem-level characterization although forceability by concepts is not proved.
& Reclassify it as a candidate obstruction and finite-network diagnostic. \\
O3 & \sevmajor & The current-status table contains both the superseded
``three quadrants plus mixed witness'' state and the later whole-fragment theorem.
& Keep one current row; move the older result to a chronology. \\
O4 & \sevmajor & Scope is easy to overread: the result is bare concept
satisfiability, not general TBoxes or unrestricted subsumption.  The raw fragment
is not negation-closed. & Put the restriction in the abstract and theorem box;
state when \(C\sqcap\neg D\) leaves the fragment. \\
O5 & \sevmajor & ``No false positive or false negative ever occurred'' conflicts
with the disclosed retracted quasimodel's false rejection and conflates labelled
regressions with oracle comparisons. & Name the current implementation and corpus;
separate 99 labelled cases from the remaining comparison runs. \\
O6 & \sevmajor & The artifact map attributes both normal-form directions to
\code{RCC5NormalForm.lean}; the final frame and set assembly is in
\code{POFreeLift.lean}. & Cite \code{odNet\_frame}, arbitrary-set interpretation, and
\code{satisfiable\_iff\_set} at their actual locations. \\
O7 & \sevmajor & The \(\Pi^0_1\)/dovetail row says ``Lean-certified,'' but the
RCC5-specific completeness and formal proof enumeration are assumptions, not an
instance. & Label only the generic race combinator source-audited and reportedly
kernel-checked. \\
O8 & \sevmod & Failure of known grid constructions is sometimes stated as ``no
concept can enforce a grid.'' & Say ``no known encoding'' unless a genuine
non-definability theorem is proved. \\
O9 & \sevmod & Patchwork is glossed as triple/path consistency implying global
consistency and rhetorically credited with resisting grids. & Separate overlap
amalgamation from path-consistency completeness and cite the exact theorem used. \\
O10 & \sevmod & ``Every load-bearing claim mechanically checked'' and ``local
algebra exhausted'' exceed the evidence. & Enumerate source-audited statements; label
F6, W2', selector forceability, and automatic-cover completeness conjectural. \\
O11 & \sevmod & Novelty and prior-art framing need tightening: abstract RCC5 set
representability predates this development; the plain concrete-domain contrast is
too categorical. & Credit relation-algebraic representation results and compare
the role-forming predicate operator separately from plain \(\mathcal{ALC}(D)\). \\
O12 & \sevminor & Review counts drift (17 versus 20), mutable \code{master} URLs
are used, the overview repeats superseded material, and the GIS example turns
membership into equality without a nominal. & Pin a release/commit, add CI/build
manifest, reconcile counts, shorten the polemical appendix, and correct the example. \\
\bottomrule
\end{longtable}
\normalsize

\subsection{The missing dependency chain}

The overview's central editorial problem is visible in one dependency chain:
\[
\begin{aligned}
\boxed{\text{semantic F6}}
&\;\mathrel{\mathop{\Longrightarrow}^{?}}\;
\boxed{\text{live-neighbourhood uniformization}}\\[-1pt]
&\;\mathrel{\mathop{\Longrightarrow}^{?}}\;
\boxed{\text{complete finite-code enumeration}}
\;\Longrightarrow\;
\boxed{\text{decision procedure}}.
\end{aligned}
\]
The last implication is formalized in conditional form.  The first two are not.
The theorem \code{decidableSat\_of\_finScheme} in
\code{Round19Transport.lean} consumes an explicit enumeration and a
\code{complete} hypothesis.  In
\code{ForcingReduction.lean}, the RCC5-specific fact is supplied through an
\code{adequacy} field, while its \code{F6fin} is qualitative existence, not an
algorithm computing a bound from the input concept.  The width-barrier prose
itself says that the coarse W2' statement is false and proposes a finer
uniformization lemma.  Calling only F6 ``the remaining theorem'' therefore
hides a theorem-level bridge.

There is also a weaker, unbounded dovetail route.  If finite certificates have
a computably enumerable syntax, their checker is decidable and sound, and every
satisfiable input has \emph{some} accepted code, then no formula-computable size
bound is needed: enumerate certificates in parallel with first-order
refutations.  The snapshot formalizes this abstract race but neither instantiates
the RCC5 certificate-completeness premise nor formalizes the refutation side.
Thus ``no computable bound'' is fatal only to the direct finite-search route;
the deeper gap in both routes is the unproved RCC5 regularization/adequacy
theorem for one exact effective certificate syntax.

There is a harmless generic sense in which decidability is equivalent to an
effective complete certificate relation.  It does not follow that this
particular live-width F6 is necessary, or that failure of this architecture
would imply undecidability.  The paper should use ``remaining gap in this
architecture,'' not ``the irreducible core of the problem.''

\subsection{Scope of the reported fragment theorem}

After normalization, the fragment forbids positive \(\forall\PO\).  On raw
syntax the polarity-correct predicate also forbids a negative
\(\exists\PO\), because normalization turns it into \(\forall\PO\).  It permits
arbitrary existential modalities, including \(\exists\PO\), and all other
universal RCC5 modalities.

This is an expressive concept-satisfiability fragment, but it is not closed
under negation and is not automatically stable under the usual reduction of
subsumption to unsatisfiability of \(C\sqcap\neg D\).  General TBox
internalization can likewise introduce \(\forall\PO\).  These limits do not
diminish the theorem; stating them prominently makes it a credible submission
candidate after the build and literature prerequisites in Section~7 are met.

\subsection{Recommended publication split}

\begin{resultbox}[Paper A: focused formal result]
Lead with the ordered-disjoint normal form, the cone control space, full-
language completeness, fragment soundness, the raw polarity theorem, and the
abstract/arbitrary-set-semantics equivalence.  Include a pinned build, axiom footprint,
complexity bound (or an honest statement that it remains loose), and a compact
comparison with prior RCC5 representation work.  This appears to be a
substantial paper after an external clean Lean build and literature sweep.
\end{resultbox}

\begin{cautionbox}[Paper B: research program and methodology]
Present F6, selector forceability, mosaics, automatic presentations, and the
history of failed routes as conjectures and experimental evidence.  This can be
valuable without claiming that the full decision problem has been reduced to
one proved keystone.  Move superseded theorem tables and the long AI-methodology
discussion here.
\end{cautionbox}

\section{Lean artifact audit}

\subsection{Canonical inventory}

The five canonical files total 51,322 lines.  The archive also contains many
historical copies, excerpts, and review bundles; those should not be allowed to
blur which files constitute the release theorem.

\begin{table}[ht]
\centering\small
\begin{tabularx}{\textwidth}{L{0.25\textwidth}rL{0.17\textwidth}Y}
\toprule
File & Lines & SHA-256 prefix & Actual role \\
\midrule
\code{POFreeLift.lean} & 45,375 & \code{5f85ff556d0d} & Current cone scheme,
fragment decision chain, raw syntax, and final arbitrary-set-semantics bridge. \\
\code{Round19Transport.lean} & 5,309 & \code{60d08f4c3adc} & Older
certificate-to-model soundness and conditional finite-scheme decider. \\
\code{SemiDecidability.lean} & 262 & \code{8a4346385619} & Generic dovetail
schema; no instantiated RCC5 certificate/refutation system. \\
\code{RCC5NormalForm.lean} & 231 & \code{438e82778910} & Compact forward
normal-form and core \(\eta\)-representation lemmas; not the whole converse assembly. \\
\code{ForcingReduction.lean} & 145 & \code{37f18d9b76f3} & Abstract reduction
under an RCC5-specific adequacy field; not a proof of adequacy or effective F6. \\
\bottomrule
\end{tabularx}
\caption{Canonical Lean sources in the submitted archive.}
\end{table}

\subsection{Proof-escape and axiom scan}

Across \code{formal/*.lean}, the static audit found:
\begin{itemize}
  \item no actual \code{sorry} or \code{admit};
  \item no project-specific \code{axiom} declaration (the sole apparent hit is
        prose inside a block comment);
  \item no \code{unsafe}, \code{opaque}, \code{partial}, \code{extern}, or
        \code{implemented\_by};
  \item 910 \code{\#print axioms} commands in the monolithic fragment file;
  \item no \code{native\_decide} use in \code{POFreeLift.lean}; 28 occurrences
        in \code{Round19Transport.lean} are diagnostics/non-vacuity examples,
        and two in \code{SemiDecidability.lean} are toy examples.
\end{itemize}
The archived transcripts report successful Lean 4.32.0/4.33.1 builds and the
standard footprint \code{propext}, \code{Classical.choice}, and
\code{Quot.sound} for the fragment capstones.  Since this review could not run
Lean, that sentence is an archived build claim, not an independently reproduced
kernel check.  The source scan nevertheless found no hidden problem-specific
assumption in the canonical cone theorem chain.

\subsection{What the current fragment chain really certifies}

\begin{longtable}{L{0.29\textwidth}L{0.13\textwidth}L{0.49\textwidth}}
\toprule
Declaration & Line & Role \\
\midrule
\endfirsthead
\toprule
Declaration & Line & Role \\
\midrule
\endhead
\code{coneScheme\_complete} & 42317 & From any full-language model/root,
obtain a surviving cone signature.  No \code{POFree} premise. \\
\code{unf\_truth} & 44006 & Truth lemma for the fresh-occurrence model;
\code{POFree} is used precisely in the \(\forall\PO\) branch near 44048. \\
\code{coneScheme\_sound} & 44093 & A survivor yields a model under the
fragment restriction. \\
\code{coneScheme\_correct\_at} & 44179 & Finite acceptance iff satisfiable for
the fragment. \\
\code{decidableSat\_cone} & 44196 & Executable NNF decision procedure. \\
\code{coneScheme\_unsat\_full} & 44217 & Sound full-language rejection only:
finite elimination implies unsatisfiable. \\
\code{nnfP\_correct} & 44452 & Raw syntax and negation-normal-form bridge. \\
\code{FPOFree} & 44530 & Polarity-sensitive fragment predicate on raw formulas. \\
\code{decidableFSat} & 44589 & Raw-formula fragment decider. \\
\code{satisfiable\_iff\_set} & 44948 & Abstract RCC5 satisfiability iff
arbitrary nonempty-set satisfiability. \\
\code{decidableSetSat} & 44976 & Set-semantics fragment decision instance. \\
\bottomrule
\end{longtable}

The clean summary is therefore:
\[
\underbrace{\text{full-language completeness}}_{\text{no false rejections}}
\quad+\quad
\underbrace{\text{fragment model extraction}}_{\text{no false acceptances}}
\quad=\quad
\text{fragment decidability}.
\]
For the full language, only the left half is present.

\subsection{What the other full-logic files do not supply}

\paragraph{Conditional finite scheme.}
\code{Round19Transport.lean} proves a credible certificate-to-model soundness
pipeline and a fixed finite checker's soundness.  Its final
\code{decidableSat\_of\_finScheme} theorem takes a \code{complete} argument.
No term inhabiting the RCC5 completeness premise is supplied.

\paragraph{Dovetailing.}
\code{SemiDecidability.lean} proves that a generic \code{DovetailScheme} whose
certificate and refutation sides are both sound and complete can be raced.
Those four properties are fields.  There is no formal first-order syntax,
Goedel coding, proof enumeration, completeness theorem, or \Logic{} instance.
Thus the mathematical co-r.e./\(\Pi^0_1\) placement may be right, while the
end-to-end certification label is not.

\paragraph{Forcing.}
\code{ForcingReduction.lean} is correct abstract order theory under the
\code{adequacy} field.  It does not build an adequacy instance for RCC5
presentations.  Its formal \code{F6fin} is qualitative, despite an
``effective'' comment: this is insufficient for the direct bounded-search
theorem, though it could support dovetailing after an actual enumerable
certificate language, checker-completeness proof, and refutation enumerator are
supplied.

\subsection{Reproduced executable evidence}

\begin{table}[ht]
\centering\small
\begin{tabularx}{\textwidth}{L{0.18\textwidth}Y L{0.20\textwidth}}
\toprule
Probe & Reproduced scope & Result \\
\midrule
\code{wp84} & 1,048,576 raw two-point and 1,953,125 functional three-point
structures; 1,050 standard-translation triples & Zero mismatches \\
\code{wp88} & 36,000 ordered-disjoint networks through size nine & All five
relations reproduced by \(\eta\); zero failures \\
\code{wp133} & 290 fresh-occurrence prefixes & Zero frame violations; reported
truth residues attributed to omitted controls \\
\code{wp134} & Seven curated elimination diagnostics & All expected outcomes \\
\code{wp135} & 2,415 fragment and 2,431 full satisfiable instances plus
mutations/boundary controls & All tested completeness obligations passed \\
\code{wp87} & 243 random/curated concepts, including depth five & Two reasoners
agreed; zero reported disagreements \\
Cover-tree stress & 911 tests (50 SAT, 13 UNSAT, 36 adversarial, 512 triples,
300 deeper random tests) & 911/911; zero mismatch \\
\bottomrule
\end{tabularx}
\caption{Finite regression evidence reproduced in this review.}
\end{table}

These runs materially increase confidence in definitions and regressions.  They
cannot establish an unbounded model theorem, and some comparisons share the
same RCC5 table or expected labels.  They should be advertised as reproducible
testing, not as independent proof.

\section{Renewed attack on full decidability}

\subsection{The ordered-disjoint core}

The normal form may be stated as a structure \((X,<,\mathrel{\#})\) where
\(<\) is a strict partial order, \(\#\) is symmetric, irreflexive, disjoint
from comparability, and downward hereditary:
\[
x<y\ \wedge\ y\# z\quad\Longrightarrow\quad x\# z,
\]
with the symmetric-coordinate version following from symmetry.  Then
\[
\begin{aligned}
\PP(x,y)&\iff x<y, &
\PPI(x,y)&\iff y<x,\\
\DR(x,y)&\iff x\#y, &
\EQ(x,y)&\iff x=y,\\
\PO(x,y)&\iff x\ne y\wedge x\not<y\wedge y\not<x\wedge\neg(x\#y).
\end{aligned}
\]
Throughout, \(X\ne\varnothing\) and strong \EQ{} means
\(\EQ(x,y)\leftrightarrow x=y\).  In the submitted set bridge, distinct points
are represented injectively by nonempty sets, with \PP{} as proper subset,
\PPI{} as proper superset, \DR{} as empty intersection, and \PO{} as nonempty
intersection with neither set containing the other.  This is arbitrary
nonempty-set semantics; it is not by itself a claim about one fixed full class
of regular-closed regions or boundary-sensitive RCC semantics.

\paragraph{Relation to the original split-forest idea.}
The ordered-disjoint normal form is best understood as the exact semantic core
of the split-forest picture, rather than as a competing construction.  Its
strict order is the vertical \PP/\PPI{} nesting axis; its hereditary conflict
relation is the horizontal \DR{} relation and precisely formalizes the
downward ``shadow'' propagation used by the split-forest argument; and \PO{} is
the remaining horizontal relation once equality, comparability, and conflict
have been excluded.  The difference is proof-theoretic strength.  An
ordered-disjoint structure may still have an arbitrary partial order, including
diamonds and several incomparable superparts.  A split forest additionally
tries to replace that order by a formula-relative tree cover, splitting join
points into occurrence copies and controlling their horizontal interfaces by
finite profiles, patchwork, or a quotient.  Thus the certified normal form
establishes the algebraic substrate, but it does not establish the bounded-width
split-cover theorem needed for the original full-logic decision route.  The
obstruction is especially sharp for \(\forall\PO\): splitting or unravelling can
create new pairs that are neither comparable nor disjoint and hence become
residual \PO-neighbours, which universal \PO{} formulas can observe.  In the
\(\forall\PO\)-free fragment those extra residual edges cannot violate a
universal \PO{} obligation, explaining why the certified fresh-occurrence
unfolding retains the split-forest shape without requiring its older quotient
and bounded-interface machinery.  For the full logic, the missing bridge is a
formula-relative split-cover theorem for ordered-disjoint structures that
preserves all \(\forall\PO\) obligations.

The full problem is therefore modal satisfiability over total five-way
labellings induced by an order with hereditary conflict.  The hard part is not
choosing one witness; it is coordinating the residual \PO{} relation among
many witnesses after order and conflict closure.

\subsection{Four boundary formulas}

The following formulas are small enough to become permanent formal regressions.
They distinguish three inadequate repairs.

\subsubsection{A direct \texorpdfstring{\(\forall\PO\)/\(\exists\PO\)}{forall-PO/exists-PO} clash}

\[
C_{\mathrm{dir}}=\forall\PO.A\ \sqcap\ \exists\PO.\neg A.
\]
This is immediately unsatisfiable: the existential witness is a \PO-neighbour
and must satisfy both \(A\) and \(\neg A\).  When the current cone control is
run with its \(\forall\PO\) information erased, the finite control survives and
accepts.  Thus simply reusing fragment extraction for the full language is
impossible.

\subsubsection{A cross-branch composition clash}

\[
C_{\mathrm{joint}}=
\exists\PP.(\forall\PPI.A\sqcap\forall\PO.A)
\ \sqcap\ \exists\PO.\neg A.
\]
Let \(y\) be the \PP-witness and \(z\) the \PO-witness from the root \(x\).
Then \(y\PPI x\) and \(x\PO z\), so the RCC5 table forces
\(\operatorname{rel}(y,z)\in\{\PPI,\PO\}\).  Either universal at \(y\)
forces \(A(z)\), contradicting \(\neg A(z)\).  Independent witness births miss
this joint obligation.

\subsubsection{A non-singleton-composition clash with no \texorpdfstring{\(\exists\PO\)}{exists-PO}}

Let
\[
\begin{aligned}
C_0={}&\exists\PPI.(\exists\PP.A)
\sqcap\forall\PP.\neg A
\sqcap\forall\PPI.\neg A
\sqcap\neg A,\\
C_{\mathrm{bad}}={}&C_0\sqcap\forall\PO.\neg A.
\end{aligned}
\]
If \(y\PP x\) and \(y\PP z\) with \(A(z)\), then
\[
\operatorname{rel}(x,z)\in
\operatorname{comp}(\PPI,\PP)=\{\EQ,\PP,\PPI,\PO\}.
\]
Under the submitted strong-\EQ{}, nonempty-region semantics, the composition
entry excludes \DR{}.  The four root constraints eliminate its four labels, so
\(C_{\mathrm{bad}}\) is unsatisfiable.  The archived diagnostic reports a
three-signature post-fixed family for the current unary cone architecture: the
\(x,z\) pair is constrained by a non-singleton composition entry and is never
emitted as an explicit \(\exists\PO\) demand.  I did not independently
reconstruct that exact family in this run, so its survival should be treated as
a source/prose diagnostic; the semantic unsatisfiability argument itself is
direct.  In any case, strengthening only explicit \PO-witness transitions does
not address this kind of hidden composition obligation.

\subsubsection{A satisfiable guard against ``make all cross-pairs \PO''}

\[
C_{\mathrm{sat}}=
\exists\PP.(A\sqcap\forall\PO.A)\ \sqcap\ \exists\PP.\neg A.
\]
Take a strict chain \(x<z<y\), use \(y\) for the first witness and \(z\) for
the second, and put \(A\) at \(y\) and \(\neg A\) at \(z\).  The witnesses
must be distinct.  At \(y\), the \(\forall\PO.A\)
condition is vacuous for these points because \(z<y\).  Thus
\(C_{\mathrm{sat}}\) is satisfiable.  Any blanket safety rule requiring
distinct independently created witnesses to be pairwise \PO-compatible would
reject it.

\begin{verdictbox}[Consequence for the current cone architecture]
A successful full construction must somehow account for relation choices and
triangle obligations among contemporaneous witnesses; an explicit joint matrix
is the most direct route.  The current endpoint types and per-demand target
tests do not do so, while a universal ``fresh pairs are PO'' default is ruled
out by \(C_{\mathrm{sat}}\).  These examples refute the submitted extraction and
the named naive repairs; they do not prove that every possible unary-state or
local-compatibility refinement must fail.
\end{verdictbox}

\subsection{Why ordinary finite types are not enough}

\paragraph{No finite cardinal bound.}
The concept
\[
C_{\infty}=\exists\PP.\top\ \sqcap\
\forall\PP.\exists\PP.\top
\]
forces an infinite ascending strict-\PP{} chain.  Therefore a finite-model
certificate cannot be the target of a model-certificate route.  Such a route
must be able to
represent infinite models finitely; this observation does not exclude a
fundamentally different decision method.

\paragraph{Filtration collapses strictness.}
Repeated points on the chain can have the same closure type.  Identifying them
turns \PP{} into equality or creates a \PP{} self-loop.  Unrolling a quotient
does not automatically help: new pairs between copies receive residual \PO,
which can violate positive \(\forall\PO\) obligations.

\paragraph{Spectra do not determine horizontal incidence.}
Even profiles enriched with lower/upper type spectra fail to record the
\DR/\PO{} matrix.  For every \(n\), take
\[
X_n=\{b\}\cup\{L_i,U_i,A_i:0\le i<n\},
\]
order only \(b<U_i\), and let conflict be the symmetric closure of
\(b\#L_i\) and \(L_i\#A_i\).  Give all members of each named family the same
unary type.  Every family then has a uniform lower/upper spectrum, while
\[
L_i\DR A_j\quad\Longleftrightarrow\quad i=j,
\qquad
L_i\PO A_j\quad\Longleftrightarrow\quad i\ne j.
\]
A quotient whose conflict choice factors only through profiles must smear the
diagonal or lose a required edge.  The example does \emph{not} prove
undecidability: a fixed concept might have a more regular alternative model.
It motivates one candidate regularization lemma---a formula-bounded
continuation index for conflict rows, or an equally effective finite
presentation---without proving that this is the unique possible route.

\newpage
\subsection{A witness-generation base lemma}

\begin{resultbox}[Witness-generated countable-model property]
If a pointed full \Logic{} concept is satisfiable, it has a countable induced
submodel generated by choosing one fixed witness for every true existential
closure formula at every retained point.  If the closure contains \(k\)
existential formulas, the chosen witness-generation graph has branching at most
\(k\).  The semantic RCC5 graph remains total, so this is not a bounded-degree
model theorem.
\end{resultbox}

\paragraph{Proof.}
Fix once and for all a selector \(w(x,E)\) choosing one original-model witness
for every pair consisting of a point \(x\) and a true existential formula \(E\)
from the NNF closure at \(x\).  With \(W_0\) the satisfying root, put
\[
W_{m+1}=W_m\cup\{w(x,E):x\in W_m,\ E\text{ is true and existential at }x\}.
\]
The union over \(\omega\) is countable, and each point has at most \(k\) outgoing
edges in this chosen generation graph.
Restrict all predicates and relations to that union.  The RCC5 frame conditions
are universal, hence inherited by an induced substructure.  Atomic literals are
unchanged; both conjuncts of a true conjunction and one disjunct of a true
disjunction remain true; universal
formulas only become easier after restriction; and each true existential has a
selected witness.  Induction over NNF formulas preserves truth at retained
points, in particular at the root. \(\square\)

This lemma is Lean-friendly and should be formalized.  It is only base camp.
A recursive finitely branching tree may have no computable path, much less an
automatic one.  Countability, compactness, and finite branching do not supply
the required regular model.

\subsection{A two-layer partial-order encoding}

The review found a concise way to absorb conflict heredity into ordinary
transitivity.  Given \((X,<,\#)\), form two copies \(Lx,Rx\) and define a
binary relation \(Q\) by
\[
\begin{array}{rclcrcl}
Lx\,Q\,Ly &\Longleftrightarrow& x<y,
&\qquad& Rx\,Q\,Ry &\Longleftrightarrow& y<x,\\
Lx\,Q\,Ry &\Longleftrightarrow& x\#y,
&& Rx\,Q\,Ly &\Longleftrightarrow& \bot.
\end{array}
\]

\paragraph{Lemma.}
\(Q\) is a strict partial order.

\paragraph{Proof sketch.}
Irreflexivity of \(Q\) follows from strictness of \(<\) on each same-sort copy;
the copies are disjoint, so a cross edge is never a \(Q\)-self-loop.
The all-left and all-right transitivity cases are order transitivity.  The mixed
case \(LxQLyQRz\) is exactly
\(x<y\wedge y\#z\Rightarrow x\#z\); the case
\(LxQRyQRz\) is the symmetric-coordinate hereditary law.  No relation leaves
the right copy for the left, so there are no other cases. \(\square\)

Let \(f\) swap \(Lx\) and \(Rx\).  On the left copy one recovers
\[
x<y\iff Q(x,y),\qquad x\#y\iff Q(x,f(y)),
\]
and hence defines \PO{} as the residual.  The original two-variable standard
translation now uses one transitive relation \(Q\), unary sort predicates, and
the matching involution \(f\).

\newpage
\begin{cautionbox}[Why this is not yet a decidability reduction]
For an exact two-copy image, the target class must enforce an involution
\(f^2=\mathrm{id}\) that swaps the left and right sorts, prohibit right-to-left
\(Q\)-edges, and require, for all left-sort \(x,y\),
\[
Q(x,y)\leftrightarrow Q(f(y),f(x)),\qquad
Q(x,f(y))\leftrightarrow Q(y,f(x)),\qquad
\neg Q(x,f(x)).
\]
The last condition is the image of conflict irreflexivity.  Standard
decidability of two-variable first-order logic with one transitive
\emph{relational} symbol does not automatically allow a unary function or an
FO-definable perfect matching.  Replacing \(f\) by its graph makes functionality
a three-variable condition in plain FO\(^2\).  Thus the construction gives a
precise literature target and perhaps a new proof route, but not a theorem that
can currently be invoked.
\end{cautionbox}

The encoding is still useful.  It packages the three-variable-looking RCC5
absorption law into transitivity and isolates synchronization, rather than
spatial composition, as the remaining issue.  A guarded/two-sorted decision
procedure specialized to this coherent involutive subclass would settle the
problem.

\subsection{Most credible positive route: automatic or mosaic covers}

An effective certificate should carry more than unary types.  Two compatible
formulations are worth pursuing.

\paragraph{Automatic presentation.}
Use a regular language of point addresses, finite-state unary type map,
synchronous relations for \PP{} and \DR{}, and regular witness relations for
each existential demand.  Define \PPI{} by converse, \EQ{} by identity, and
\PO{} residually.  Given such a finite presentation, automata operations can
check strictness, transitivity, conflict symmetry/irreflexivity, heredity,
endpoint modal constraints, and total witness service.  Soundness can then feed
the reportedly kernel-checked set representation.

\paragraph{Pair-matrix mosaics.}
A finite state should contain a complete RCC5 matrix among one representative
for each simultaneously live demand, with interfaces recording continuation
requirements.  A greatest fixed point must prove both directions:
\begin{enumerate}
  \item mosaics realized in a model survive elimination; and
  \item every surviving compatible mosaic family unfolds to one coherent
        ordered-disjoint model satisfying all \(\forall\PO\) obligations.
\end{enumerate}
The second statement is where bounded joint width or finite continuation index
reappears.  Calling it ``F6'' does not prove it.

The sharp qualitative open statement is therefore
\[
\begin{aligned}
\mathsf{EnumerableCover}(C):\quad
C\text{ satisfiable}\quad\Longrightarrow\quad
&\text{some finite automatic (or mosaic) code}\\[-2pt]
&\text{passes the fixed effective checker}.
\end{aligned}
\]
Together with a formal enumerable refutation side, this is enough for
decidability by dovetailing.  A stronger computable bound on code size would
instead give a direct finite search.  No proof of either cover statement is
present.  A useful intermediate fragment would
bound the number of distinct \DR{} rows between equal closure profiles and show
that the bound is preserved by witness generation.

\section{Literature and novelty audit}

\subsection{What the primary comparator establishes}

Lutz and Wolter's 2006 paper is the closest primary comparator
\cite{LutzWolter2006}.  The archived PDF confirms the exact citation as
LMCS 2(2:5), not 2(2:2).  Their representation theorems identify general RCC5
structures with the relevant substructure semantics; their undecidability
theorem applies to supremum-closed/full region classes, not to the general
substructure class matching the submitted target.  They prove recursive
enumerability of the corresponding validity logic and explicitly leave the
general RCC5 case open.  Consequently, their result does not supply either
full decidability or undecidability here.

The co-r.e./\(\Pi^0_1\) satisfiability upper placement used in the overview is
consistent with this: validity is r.e., so unsatisfiability of a concept can be
enumerated after the standard translation.  What is missing from Lean is the
actual formal proof system/encoding and its completeness, not the abstract idea
of dovetailing two semideciders.

\subsection{Representation and concrete-domain context}

Abstract/general RCC5 set representability is not new as a bare mathematical
fact.  The relation-algebraic literature of Duentsch, Wang, and McCloskey
\cite{DuentschWangMcCloskey2001}, Li and Ying's analysis of models and the
composition table \cite{LiYing2003}, and Lutz--Wolter's representation results
must be credited.  The potentially novel contribution here is the particular
ordered-disjoint formulation, its constructive \(\eta\) representation, and
the reportedly kernel-checked end-to-end use in a description-logic fragment.

The overview should also narrow its claim that concrete-domain formalisms
cannot form relations.  Plain \(\mathcal{ALC}(D)\) is indeed feature-mediated
and two-sorted, but Haarslev, Lutz, and Moeller explicitly introduced a
role-forming predicate operator over feature chains
\cite{HaarslevLutzMoeller1999}.  That formalism is still not the primitive,
single-sorted, total RCC5 role partition studied here; it is simply a closer
comparator than the overview currently admits.

\subsection{Event structures and the doubled-poset target}

The ordered-disjoint frame has the same order/conflict signature and a dual
heredity law as an event-structure skeleton: event causality corresponds to the
\emph{reverse} of \(<\,{=}\,\PP\), hereditary \DR{} is conflict, and residual
\PO{} is concurrency \cite{NielsenPlotkinWinskel1981,Winskel1987}.  An arbitrary
frame here need not be a prime event structure: classical definitions commonly
require finite causal past, which is absent from the submitted frame axioms.
No direct decision theorem for this event-to-event modal
signature was located in the inspected corpus; standard event-structure logics
often reason over configurations and enabling transitions instead.  The
connection nevertheless supplies mature language for unfoldings, regular event
structures, and finite-state presentations, and deserves a precise comparison.

Szwast and Tendera prove decidability of function-free relational FO\(^2\) with
one designated transitive relation and auxiliary relational predicates
\cite{SzwastTendera2013}.  The theorem is a comparator, not a corollary: the
direct \((<,\#)\) theory needs the three-variable hereditary axiom, while the
two-layer encoding needs a synchronized unary involution/perfect matching.
This exact enriched signature is the right literature question.

\subsection{Search limitation}

The supplied archive contains the primary 2006 paper and a substantial related-
work record.  No inspected primary source in that corpus settles full
\Logic{} satisfiability.  Live web/database searching was attempted but was not
available from this run, so this report does \emph{not} certify the stronger
bibliographic statement ``still open as of September 2026.''  Before a novelty
or priority claim is published, run and record searches in DBLP, MathSciNet,
zbMATH, Scopus/Web of Science, Google Scholar, and citation graphs for the 2006
paper, including terms for RCC5 modal logic, event structures, one transitive
relation with functions, and regular/automatic models.

\section{Recommendations and research plan}

\subsection{Priority actions}

\begin{longtable}{L{0.10\textwidth}L{0.29\textwidth}L{0.51\textwidth}}
\toprule
Priority & Action & Acceptance criterion \\
\midrule
\endfirsthead
\toprule
Priority & Action & Acceptance criterion \\
\midrule
\endhead
P0 & Freeze a reproducible Lean release & Add \code{lean-toolchain}, a minimal
Lake project, shared base module, and CI that builds the five canonical targets,
rejects \code{sorry}/custom axioms, and records expected axiom footprints. \\
P0 & Correct the overview's theorem status & One current fragment theorem;
F6/uniformization/checker arrows shown separately; no iff selector claim; bare
concept/TBox/subsumption scope explicit. \\
P0 & Add full-boundary regressions & Formalize \(C_{\mathrm{dir}}\),
\(C_{\mathrm{joint}}\), \(C_{\mathrm{bad}}\), and \(C_{\mathrm{sat}}\), including
survival/failure statements for candidate controls. \\
P1 & Extract a trusted shared core & One imported definition of atoms,
composition, syntax, semantics, and ordered-disjoint representation; eliminate
large copied tables across canonical files. \\
P1 & Formalize countable witness generation & A carrier-polymorphic theorem
constructing the induced countable submodel and proving NNF truth preservation. \\
P1 & Build certificate soundness first & Define an automatic or pair-mosaic
certificate in finite syntax; prove its checker decidable and its accepted
objects yield models, without hiding completeness in a structure field. \\
P1 & State the hard hypothesis honestly & A named
\(\mathsf{EnumerableCover}\) theorem for one effective finite-code grammar;
separately label any stronger formula-computable bound used for direct search. \\
P2 & Attack bounded-row fragments & Prove completeness when each closure-profile
pair has at most \(m\) distinct hereditary-conflict rows; determine whether
\(m\) can be bounded from the formula. \\
P2 & Develop the doubled-poset route & Formalize the two-copy equivalence and
survey/derive decidability for FO\(^2\) over one strict order with the required
coherent involution. \\
P2 & Run a reproducible literature review & Publish queries, dates, databases,
inclusion criteria, and a citation graph; have a human domain expert verify the
priority conclusion. \\
\bottomrule
\end{longtable}

\subsection{Go/no-go milestones for full decidability}

\begin{enumerate}
  \item \textbf{Semantic regression gate.}  Any new local rule must reject
        \(C_{\mathrm{bad}}\) and accept \(C_{\mathrm{sat}}\).  Failure means the
        architecture is not ready for a completeness proof.
  \item \textbf{Finite-code gate.}  Certificates must be inductive finite data
        with a terminating checker.  Higher-order functions or an unbounded
        semantic labelling are not decision certificates.
  \item \textbf{Soundness gate.}  Prove checker acceptance produces an
        ordered-disjoint model satisfying the formula, with no completeness
        field in the soundness statement.
  \item \textbf{Regularization gate.}  Prove that every satisfiable formula has
        an accepted finite code in that exact effective syntax (or a stronger
        formula-computable bound if direct finite search is claimed).  This is
        the decisive theorem for this positive route.  If it fails, report the countermodel
        pattern rather than shifting the premise to another abstraction.
  \item \textbf{End-to-end gate.}  Instantiate the executable decision theorem,
        print its axioms, run it on all boundary formulas, and reproduce it from
        a clean tagged checkout in CI.
\end{enumerate}

\subsection{What not to infer}

Three negative conclusions would be premature:
\begin{itemize}
  \item Failure of the unary cone scheme does not imply undecidability.
  \item Selector matrices in arbitrary finite RCC5 frames do not show that one
        fixed concept forces unbounded selectors.
  \item A countable finitely branching witness construction does not imply an
        automatic or computable model.
\end{itemize}
Conversely, the absence of a successful grid reduction does not imply
decidability.  The honest frontier is a regularization theorem or a genuinely
faithful undecidability encoding.

\section{Final assessment}

\begin{verdictbox}[Final answer]
The attached work has crossed an important line: the \(\forall\PO\)-free
fragment is no longer merely an experimental conjecture but is backed by a
detailed, apparently escape-free Lean theorem chain and substantial regression
evidence.  That is the result to develop into a focused submission after a
pinned independent build and current novelty review.

The full \Logic{} decision problem is unresolved in this snapshot.  The
submitted cone extraction is not repaired merely by erasing or adding the
endpoint-only \(\forall\PO\) checks considered here; the formulas in
Section~5.2 show why.  The most credible positive target is an enumerable
joint-pair/mosaic or automatic-presentation cover theorem (with a computable
bound if direct finite search is desired).  The two-layer partial-
order encoding offers a second, sharply stated route, but the required matching
involution is a nontrivial enrichment not covered by the standard FO\(^2\)
theorem.

Publication verdict on the overview: \textbf{major rewrite / reject in current
formal-results form}.  Publication verdict on a separated fragment/formalization
paper after a pinned independent Lean build: \textbf{promising and likely
publishable, subject also to a current novelty review}.  Research verdict on the
full problem: \textbf{unresolved in the reviewed snapshot and corpus; a
sufficient global regularization theorem for the leading positive route is now
more precisely stated, but not proved}.
\end{verdictbox}

\appendix

\section{Artifact hashes and reproduction notes}\label{app:repro}

\small
\begin{longtable}{L{0.32\textwidth}L{0.61\textwidth}}
\toprule
Artifact & SHA-256 \\
\midrule
\endfirsthead
\toprule
Artifact & SHA-256 \\
\midrule
\endhead
\code{alcircc5-master.zip} & \nolinkurl{4a0415b3e3202595d1b1290e6cc741068d0423c0c2877e815b9b5fcbe96bd2a8} \\
\code{papers/overview\_arxiv.tex} & \nolinkurl{be28d3dace3a1cf0bacbaaa253872c38e76c32a9995d3e21c70140b8e656942e} \\
\code{papers/overview\_arxiv.pdf} & \nolinkurl{e684c506ac501219a93cf1b733d52dfcaa3a5c63f983bf6a651199000440322a} \\
\code{formal/POFreeLift.lean} & \nolinkurl{5f85ff556d0d67205432311f136a90335b4a0540576604ec0f395e4aa3db28ae} \\
\code{formal/Round19Transport.lean} & \nolinkurl{60d08f4c3adcd44b90820d764911750cc3e19072a7e84e6e87174e4a36c82678} \\
\code{formal/SemiDecidability.lean} & \nolinkurl{8a43463856193ca6ad053498502c17239fff8c9d33d7aab594e4f5e9d1077347} \\
\code{formal/RCC5NormalForm.lean} & \nolinkurl{438e8277891054c82050d2d1c83065100881b0630692b72c7c71a5ee80d127ad} \\
\code{formal/ForcingReduction.lean} & \nolinkurl{37f18d9b76f38461dd310a44adc768216af0070a88bda0dd482ce761269b3351} \\
\bottomrule
\end{longtable}
\normalsize

The following classes of checks were run from the extracted archive:
\begin{itemize}
  \item exact hash, line-count, and forbidden-token scans over canonical Lean;
  \item theorem-location/dependency inspection around the cone capstones;
  \item \code{wp84\_fo\_pi01\_transcription.py} and
        \code{wp88\_canonical\_representation.py};
  \item \code{wp133}, \code{wp134}, \code{wp135}, and \code{wp87};
  \item \code{src/stress\_test\_cover\_tree.py};
  \item direct execution of the full-logic boundary script showing acceptance
        of \(C_{\mathrm{dir}}\) and \(C_{\mathrm{joint}}\) when
        \(\forall\PO\) obligations are erased.
\end{itemize}
No fresh Lean compilation was possible for the reasons stated in Section~1.1.

\section{Claim-to-evidence matrix}

\begin{longtable}{L{0.34\textwidth}L{0.24\textwidth}L{0.34\textwidth}}
\toprule
Claim & Best evidence in snapshot & What remains \\
\midrule
\endfirsthead
\toprule
Claim & Best evidence in snapshot & What remains \\
\midrule
\endhead
Ordered-disjoint frames yield RCC5 models & Lean normal-form and set lemmas;
36,000-network probe & Clean modular rebuild and corrected theorem map \\
Full models yield surviving cone signatures & \code{coneScheme\_complete} &
Independent pinned build \\
Surviving fragment signature yields a model & \code{unf\_truth},
\code{coneScheme\_sound} & Independent pinned build \\
Raw \(\forall\PO\)-free SAT is decidable & \code{FPOFree},
\code{decidableFSat} & Independent pinned build, complexity statement, and
release engineering \\
Concrete set SAT matches abstract SAT & \code{satisfiable\_iff\_set} & Credit
prior representation theorems precisely and reproduce a pinned build \\
Full SAT is decidable from a fixed complete finite scheme & Conditional Lean
combinator & Actual bounded scheme and proof of completeness \\
SAT is in \(\Pi^0_1\) & First-order transcription plus Lutz--Wolter r.e.
validity argument & Formal encoding/proof calculus instance in Lean \\
F6 yields complete finite certificates & Prose sketches / abstract adequacy
field & RCC5 adequacy/uniformization into an enumerable finite-code grammar;
optional bound for direct search \\
Full cone soundness & None; boundary formulas expose the submitted gap and
defeat named naive repairs & Joint pair/mosaic or different model construction \\
No solution is established in the reviewed snapshot & 2006 primary source plus
supplied corpus & Fresh reproducible 2026 literature review \\
\bottomrule
\end{longtable}

\begin{thebibliography}{99}

\bibitem{LutzWolter2006}
C. Lutz and F. Wolter.
\newblock Modal logics of topological relations.
\newblock \emph{Logical Methods in Computer Science}, 2(2:5):1--41, 2006.
\newblock \url{https://doi.org/10.2168/LMCS-2(2:5)2006}.

\bibitem{DuentschWangMcCloskey2001}
I. Duentsch, H. Wang, and S. McCloskey.
\newblock A relation-algebraic approach to the Region Connection Calculus.
\newblock \emph{Theoretical Computer Science}, 255(1--2):63--83, 2001.
\newblock \url{https://doi.org/10.1016/S0304-3975(99)00123-5}.

\bibitem{LiYing2003}
S. Li and M. Ying.
\newblock Region Connection Calculus: its models and composition table.
\newblock \emph{Artificial Intelligence}, 145:121--146, 2003.
\newblock \url{https://doi.org/10.1016/S0004-3702(02)00372-7}.

\bibitem{HaarslevLutzMoeller1999}
V. Haarslev, C. Lutz, and R. Moeller.
\newblock A description logic with concrete domains and a role-forming
predicate operator.
\newblock \emph{Journal of Logic and Computation}, 9(3):351--384, 1999.
\newblock \url{https://doi.org/10.1093/logcom/9.3.351}.

\bibitem{NielsenPlotkinWinskel1981}
M. Nielsen, G. Plotkin, and G. Winskel.
\newblock Petri nets, event structures and domains, Part I.
\newblock \emph{Theoretical Computer Science}, 13(1):85--108, 1981.
\newblock \url{https://doi.org/10.1016/0304-3975(81)90112-2}.

\bibitem{Winskel1987}
G. Winskel.
\newblock Event structures.
\newblock In \emph{Petri Nets: Applications and Relationships to Other Models
of Concurrency}, LNCS 255, pages 325--392, 1987.
\newblock \url{https://doi.org/10.1007/3-540-17906-2_31}.

\bibitem{SzwastTendera2013}
W. Szwast and L. Tendera.
\newblock FO\({}^2\) with one transitive relation is decidable.
\newblock In \emph{STACS 2013}, LIPIcs 20, pages 317--328, 2013.
\newblock \url{https://doi.org/10.4230/LIPIcs.STACS.2013.317}.

\bibitem{Lean4}
L. de Moura and S. Ullrich.
\newblock The Lean 4 theorem prover and programming language.
\newblock In \emph{CADE 28}, LNCS 12699, pages 625--635, 2021.
\newblock \url{https://doi.org/10.1007/978-3-030-79876-5_37}.

\bibitem{Archive2026}
M. Wessel et al.
\newblock ALCI-RCC5 research repository, archived commit
\nolinkurl{795b2707b72e927c56b465ef0413b77aeb9a66a6}, 2026.
\newblock \url{https://github.com/lambdamikel/alcircc5/tree/795b2707b72e927c56b465ef0413b77aeb9a66a6}.

\end{thebibliography}

\end{document}
